% FILE: ~/GPLC/MEMORIA/sections/validacion.tex

\newpage
\section{Validación}



El fin de este capítulo es presentar un conjunto de casos que cumplen
un doble propósito: \textbf{1.} aportar evidencia de que el intérprete
implementa correctamente el lenguaje GPLC y \textbf{2.}  ilustrar el
uso práctico del lenguaje mediante programas pequeños pero
representativos.  La evidencia obtenida se basa en el contraste entre
los resultados producidos por el intérprete y el comportamiento
esperado: en los casos exitosos, las distribuciones retornadas y sus
valores con ascripcion, y en los casos fallidos, los diagnósticos de
error reportados.

A diferencia de una verificación formal completa, la validación que se
presenta aquí es funcional y dirigida por casos.  Cada programa está
diseñado para cubrir una expresión específica (o una interacción
acotada entre expresiones) y contrastar el comportamiento observado
con las reglas y decisiones de diseño descritas en capítulos previos.
Este enfoque resulta adecuado para GPLC, ya que lo que se busca
contrastar se refleja en hechos directamente observables: la
aceptación de programas bien formados y la producción de
distribuciones con valores y tipos coherentes; la propagación de
probabilidades concretas y desconocidas de acuerdo con las reglas del
lenguaje; y el rechazo estático o dinámico de programas inválidos
mediante diagnósticos precisos con localización y explicación de la
inconsistencia detectada.

  Finalmente, estos casos también cumplen un rol ilustrativo, ya que muestran
  patrones de uso propuestos, ejemplos mínimos y casos de borde que
  aclaran qué se acepta, qué se rechaza y por qué. Por ello, el
  capítulo está organizado en torno a familias
  de expresiones, agrupando los programas por la construcción que se desea validar.


\subsection{Metodología}

Para cada caso se documentan tres elementos: \textbf{1.} el código
fuente GPLC, \textbf{2.} la salida del intérprete (ya sea una
distribución simbólica o un mensaje de error), y \textbf{3.} una
explicación breve del aspecto que se está validando. En los casos
exitosos, la validación se apoya en observar que: \textbf{i.} el
programa es aceptado por el chequeo estático, \textbf{ii.} la
distribución retornada contiene los valores esperados, y \textbf{iii.}
los tipos que figuran en cada valor son los correctos. En los casos
fallidos, se espera un rechazo estático o dinámico según corresponda,
y se valida que el intérprete reporte el error con un mensaje
informativo y una localización adecuada.

Los ejemplos se seleccionan además para cubrir el \emph{carácter
  gradual} del lenguaje: probabilidades desconocidas \verb|?|, el tipo
desconocido \verb|?|, y el efecto de estos elementos en consistencia,
ascripción y construcción de distribuciones. Con esto, el capítulo
busca evidenciar no solo que el intérprete evalúa programas, sino
también preserva las restricciones semánticas y de integridad que
motivan el diseño de GPLC para escenarios con incertidumbre e
información parcial.



\subsection{Casos}

Los casos se organizan por construcción sintáctica: selección
probabilística para distribuir probabilidad entre dos ramas,
\lstinline|let| para definir variables, lambdas para abstracción y
aplicación, \lstinline|if| para control y ascripción para imponer
tipos. Cada set de casos reúne programas mínimos que aíslan el
comportamiento bajo estudio, y luego incorpora ejemplos donde la
interacción entre componentes graduales (tipos y probabilidades
desconocidas) es relevante. Cuando corresponde, también se incluyen
casos de rechazo estático y dinámico para ilustrar las condiciones
bajo las cuales el intérprete protege la integridad semántica de
probabilidades y distribuciones.

\subsubsection{Selección Binaria Probabilística}
El código \ref{lst:choice1} muestra un caso elemental
de selección probabilística: la selección de un
número real con probabilidad concreta 0.5. El
output muestra una distribución simbólica donde
los elementos 1 y 0 están ascritos por el tipo
\lstinline|Real|, y ponderados por probabilidades simbólicas.


\begin{gplcblock}[caption={Seleción de números reales},label={lst:choice1}]
0.5 $ 1 : 0

>>>

{
	 [Real] 1. :: Real : SymProb19 ,
	 [Real] 0. :: Real : SymProb20 
}   
\end{gplcblock}




El programa \ref{lst:choice2} muestra un caso de selección
con ramas no homogéneas: valores de tipo \lstinline|Real| y
\lstinline|Bool| son seleccionados, lo cual
constituye una decisión
de diseño importante: Las ramas no
necesitan ser de un mismo tipo. Además la selección
es hecha con una probabilidad incógnita. La variable
simbólica libre asociada a dicha probabilidad queda
internamente codificada en la fórmula de la distribución
retornada.







  
\begin{gplcblock}[caption={Selección de booleano y número},label={lst:choice2}]
? $ true : 10

>>>

{
	 [Bool] true :: Bool : SymProb19 ,
	 [Real] 10. :: Real : SymProb20 
}  
\end{gplcblock}




El programa \ref{lst:choice3} muestra un ejemplo de selección anidada,
es decir, una selección con una rama que, a su vez,
es una selección. En el output se observa como los
tres valores conforman el soporte de la distribución
retornada.

\begin{gplcblock}[caption={Selección anidada},label={lst:choice3}]
0.4 $ (? $ true : false) : 10

>>>

{
	 [Bool] true :: Bool : SymProb71 ,
	 [Bool] false :: Bool : SymProb72 ,
	 [Real] 10. :: Real : SymProb73 
}  
\end{gplcblock}



En el programa \ref{lst:choice4}, la expresión de selección
tiene una ascripción de tipo de distribución, el cual
verifica estáticamente su consistencia con el tipo
inferido de la expresión. En la distribución ascrita,
el tipo \lstinline|Real| concentra el peso total de los elementos
reales, lo cual comprueba que en la relación de
consistencia la distribución de los pesos de un tipo
no afecta la consistencia entre las distribuciones. Por
lo tanto, el resultado de un juicio de consistencia está
determinado estrictamente por relaciones semánticas entre
los tipos, sin importar la sintaxis específica. Esto es
cierto en general para todos las relaciones y operaciones
entre distribuciones.



\begin{gplcblock}[caption={Selección con ascripción},label={lst:choice4}]
(0.5 $ 1 : (0.4 $ 0 : false)) :: {Real : 0.7 , Bool : 0.3}

>>>

{
	 [Real] 1. :: Real : SymProb46 ,
	 [Real] 0. :: Real : SymProb47 ,
	 [Bool] false :: Bool : SymProb48 
}  
\end{gplcblock}




  
Con un carácter más gradual, los programas \ref{lst:choice5.1}, \ref{lst:choice5.2} y \ref{lst:choice5.3}
muestran la flexibilidad que se logra con el tipo
desconocido \lstinline|?| y la probabilidad desconocida \lstinline|?|.
El programa \ref{lst:choice5.1} muestra que el tipo \lstinline|?|
--consistente con los tipos \lstinline|Real| y \lstinline|Bool|--
puede distribuir su peso entre \lstinline|Real| y \lstinline|Bool|
para lograr la consistencia con el tipo de la
expresión. El programa \ref{lst:choice5.2} evidencia que
la distribución con probabilidades desconocidas es
consistente con toda expresión cuyo tipo sea una
distribución de \lstinline|Real| y \lstinline|Bool|, o incluso una
distribución con solo uno de esos dos tipos, ya que
la probabilidad \lstinline|?| es compatible con 0.0, es decir,
puede representar también la probabilidad de un tipo
con peso nulo. Finalmente, el programa \ref{lst:choice5.3}
demuestra que se puede anotar la selección con cualquier
distribución de \lstinline|Real| y \lstinline|Bool|, dado que la selección
con probabilidad desconocida es consistente con
distribución de pesos en tipo de ascripción.


\begin{gplcblock}[caption={Distribución con tipo no especificado},label={lst:choice5.1}]
(0.3 $ true : (-1)) :: {Bool : 0.2, ? : 0.8}

>>>

{
	 [Bool] true :: Bool : SymProb21 ,
	 [Bool] true :: ? : SymProb22 ,
	 [Real] -1. :: ? : SymProb23 
}  
\end{gplcblock}

\begin{gplcblock}[caption={Distribución con probabilidades no especificadas},label={lst:choice5.2}]
(0.3 $ true : (-1)) :: {Real : ?, Bool : ?}

>>>

{
	 [Bool] true :: Bool : SymProb21 ,
	 [Real] -1. :: Real : SymProb22 
}  
\end{gplcblock}

\begin{gplcblock}[caption={Ascripción de selección con prob. \lstinline|?|},label={lst:choice5.3}]
(? $ true : (-1)) :: {Real : 0.5 , Bool : 0.5}

>>>

{
	 [Bool] true :: Bool : SymProb21 ,
	 [Real] -1. :: Real : SymProb22 
}  
\end{gplcblock}



Para demostrar con claridad el funcionamiento
efectivo del mecanismo estático de tipado, se analizan
un par de casos donde la ascripción impuesta a la
selección gatilla un error estático.

En el programa \ref{lst:choice6}, la ponderación exacta
que tiene el tipo \lstinline|Bool| en la distribución inducida
por la selección, es exactamente \lstinline|0.4|. Sin embargo,
en la ascripción, el tipo `Bool` tiene una ponderación
mál alta.


\begin{gplcblock}[caption={Inconsistencia por prob. de \lstinline|Bool|},label={lst:choice6}]
(0.4 $ false : lambda x:?.x) :: {Bool:0.45 , ?:?}

>>>

Static Error [line 1, cols 1-49]: In distribution ascription, type 0.4 * ({Bool:1}) + 0.6 * ({? -> {?:1}:1}) of term inconsistent with ascription type {Bool:0.45, ?:?}
\end{gplcblock}

Para terminar, el programa \ref{lst:choice7} falla porque
en la anotación no hay ningún tipo consistente con
el tipo funcional \lstinline|Real -> {Real:1}| que represente
su peso \lstinline|0.5| en el tipo de distribución ascrito.




\begin{gplcblock}[caption={Inconsistencia por falta de tipo funcional},label={lst:choice7}]
(0.5 $ lambda x:Real.1 : false) :: {Bool : ? , Real : ?}

>>>

Static Error [line 1, cols 1-56]: In distribution ascription, type 0.5 * ({Real -> {Real:1}:1}) + 0.5 * ({Bool:1}) of term inconsistent with ascription type {Bool:?, Real:?}
\end{gplcblock}

\subsubsection{Let}
La semántica de la operación \lstinline|let| consiste en
asignar a la variable \lstinline|x| cada uno de los
elementos contenidos en la distribución implícita
de la expresión nombrada, y luego establecer una
distribución de probabilidad que contenga las
instancias del cuerpo con cada versión de
variable \lstinline|x|. La probabilidad de cada valor asignado
a \lstinline|x|, dentro de la distribución ligada, se toma como
el peso del cuerpo en la distribución resultante.

En el primer ejemplo \ref{lst:let1}, el programa es
semánticamente equivalente a la expresión
\lstinline|(0.3 $ true : 1)|. La distribución de valores
retornada muestra la combinación probabilística
efectiva del cuerpo del let (simplemente la variable x),
evaluado en todo el rango de la variable.

La variable funciona como en el común de los lenguajes
cuando la expresión nombrada representa solo un elemento
como en los casos \lstinline|y = false| o \lstinline|id = lambdax:?.x|.

\begin{gplcblock}[caption={Definición de variable local},label={lst:let1}]
let x = (0.3 $ true : 1) in x

>>>

{
	 [Bool] true :: Bool : SymProb41 ,
	 [Real] 1. :: Real : SymProb42 
}  
\end{gplcblock}

El mayor potencial del bloque \lstinline|let| consiste
naturalmente en la operabilidad  que le otorga
a los elementos de una variable de distribución.
En el programa \ref{lst:let2.1}, la variable b es el
argumento de una operación lógica, y en el
programa \ref{lst:let2.2} la variable r participa
en una operación aritmética. Aunque no sea
visible en los pesos simbólicos, la variable
propaga a la distribución resultante el peso
probabilístico de los valores individuales
asignados a ella.


\begin{gplcblock}[caption={Uso de variable en operación lógica},label={lst:let2.1}]
let b = (0.9 $ true : false) in (true or b)

>>>

{
	 [Bool] true :: Bool : SymProb73 ,
	 [Bool] true :: Bool : SymProb74 
}  
\end{gplcblock}



\begin{gplcblock}[caption={Uso de variable en operación aritmética},label={lst:let2.2}]
let r = (? $ 0 : 10) in (7 * r)

>>>

{
	 [Real] 0. :: Real : SymProb73 ,
	 [Real] 70. :: Real : SymProb74 
}  
\end{gplcblock}

El programa \ref{lst:let3} muestra bloques \lstinline|let| anidados.
Dada la naturaleza semántica de la expresión \lstinline|let|,
la expresión final en términos de \lstinline|n| y \lstinline|b| combina
todos los valores posibles de las variables, cuatro
combinaciones en total. Además, el programa se
verifica con la ascripción de un tipo de distribución
consistente con la equiprobabilidad de \lstinline|Real|
y \lstinline|Bool|.




\begin{gplcblock}[caption={Bloques \lstinline|let| anidados},label={lst:let3}]
(let n = (? $ 1 : 7) in
let b = (0.3 $ true : false) in
(0.5 $ 21 / n : b and false))  :: {Real:0.5, Bool:0.5}

>>>

{
	 [Real] 21. :: Real : SymProb403 ,
	 [Bool] false :: Bool : SymProb404 ,
	 [Real] 21. :: Real : SymProb405 ,
	 [Bool] false :: Bool : SymProb406 ,
	 [Real] 3. :: Real : SymProb407 ,
	 [Bool] false :: Bool : SymProb408 ,
	 [Real] 3. :: Real : SymProb409 ,
	 [Bool] false :: Bool : SymProb410 
}  
\end{gplcblock}

El ejemplo \ref{lst:let4} exhibe elementos importantes de
la expresión \lstinline|let| en conexión con las expresiones
\lstinline|lambda|. es muy conveniente la asignación de una lambda
a una variable, pues permite una aplicación mucho
menos verbosa y más práctica que la aplicación de
una lambda literal directamente al argumento. En
segundo lugar, el programa demuestra la capacidad del
bloque \lstinline|let| para extraer funciones de una seleción.
La ascripción verifica que los resultados son
exclusivamente de tipo \lstinline|Real|.


\begin{gplcblock}[caption={Asignación de \lstinline|lambda| a variable},label={lst:let4}]
(let step_forward = lambdax.x+1 in
let step_back = lambday.y-1 in
let f = (? $ step_forward : step_back) in
(f 0)) :: {Real:1.0}

>>>

{
	 [Real] 1. :: Real : SymProb145 ,
	 [Real] -1. :: Real : SymProb146 
}  
\end{gplcblock}

En el contexto de los bloques \lstinline|let| y las expresiones
lambda, es pertinente mencionar que si se usa una variable que no ha
sido definida con los medios que la sintaxis ofrece, como la
asignación de una expresión en un bloque \lstinline|let|, o como
parámetro formal de una expresión lambda, entonces el intérprete
gatilla un error cuyo mensaje especifica cuál es la variable no
definida, y la ubicación exacta del código en que fue utilizada.  En
el código \ref{lst:invalid-var} se muestra el uso inválido de la
variable \lstinline|x| en el cuerpo de la expresión lambda.

\begin{gplcblock}[caption={Uso de variable no definida},label={lst:invalid-var}]
lambda y : Bool . (x + y)

>>>

Static Error [line 1, cols 19-20]: variable x not found
\end{gplcblock}


\subsubsection{Lambda}

Las expresiones lambda tienen la sintaxis \lstinline|lambda<var>:<simple>.<expr>|.
La sintaxis\lstinline| lambda<var>.<expr>| es azúcar sintático de
\lstinline|lambda<var>:?.<expr>|. Si bien el argumento de una función
debe ser una expresión de tipo simple, es decir, un valor
(número, booleano, variable, o lambda), el cuerpo
de la función puede ser cualquier clase de expresión.

En el programa del ejemplo \ref{lst:lambda1.1}, la función
\lstinline|seven_untyped| tiene un parámetro formal de tipo \lstinline|?|.
Por lo tanto, la función es capaz de recibir un argumento
con cualquier tipo simple, sin levantar un error estático
inconsistencia, porque \lstinline|?| es consistente con todo tipo
simple. (En el caso particular de esta función, tampoco
sería posible tener un error dinámico, porque el cuerpo
no tiene la variable como subexpresión). Los argumentos
1 y true es pasan exitosamente y el programa termina.

\begin{gplcblock}[caption={Función no tipada},label={lst:lambda1.1}]
let seven_untyped = lambdax:?.7 in
let a1 = (seven_untyped 1) in
let a2 = (seven_untyped true) in
a1 * a2 

>>>

{
	 [Real] 49. :: Real : SymProb85 
}  
\end{gplcblock}

En cambio, en el programa \ref{lst:lambda1.2}, la lambda
\lstinline|seven_typed| se define con un parámetro de tipo
\lstinline|Real|, explícitamente declarado. En este segundo
ejemplo, el tipado gradual rechaza estáticamente
la expresión \lstinline|(sevent_typed false)|, debido a la
inconsistencia del tipo del argumento (\lstinline|Bool|) con
el tipo del parámetro (\lstinline|Real|).

Por lo tanto, el sistema gradual prueba en esta
implementación tener la flexibilidad para recibir
argumentos de todo tipo, cuando el tipo del parámetro
se omite o se declara como \lstinline|?|, y decide en tiempo
de ejecución la compatibilidad del argumento. Por
otro lado, el grado de refinamiento del tipo declarado
del parámetro puede resultar en un rechazo estático del
programa.




\begin{gplcblock}[caption={Función explícitamente tipada},label={lst:lambda1.2}]
let seven_typed = lambdax:Real.7 in
let b1 = (seven_typed false) in
let b2 = (seven_typed 0) in
b1 + b2  

>>>

Static Error [line 2, cols 9-28]: In function, expected argument of type Real, but found Bool
\end{gplcblock}


En el lenguaje GPLC, las funciones son de primera clase,
se decir, son valores que puede a su vez pasarse como
como argumento a otra función. En el ejemplo \ref{lst:lambda2},
se define un aplicador como una función que recibe otra
función como argumento, y la aplica al número \lstinline|5|. Luego,
se define \lstinline|sqr| como función a ser aplicada. En la expresión
principal, el aplicador toma \lstinline|sqr| como función argumento
y la aplica al \lstinline|5|.

\begin{gplcblock}[caption={Función como argumento de otra función},label={lst:lambda2}]
let apply = lambda f : Real -> {Real:1} . (f 5) in
let sqr = lambda r . (r * r) in
(apply sqr)

>>>

{
	 [Real] 25. :: Real : SymProb72 
}  

\end{gplcblock}





En el programa \ref{lst:lambda3} se desmuestra cómo una función
puede operar sobre una variable ligada a los elementos
de una distribución inducida por una selección probabilística.
En la evaluación del programa, el \lstinline|64| es dividido por \lstinline|2|,\lstinline|4|, y \lstinline|8|,
distribuyendo  el peso total de una selección numérica (0.8) de
acuerdo al peso de cada valor de \lstinline|n|.

\begin{gplcblock}[caption={Función aplicada a identificador de distribución},label={lst:lambda3}]
let n = 0.5 $ ( ? $ 2 : 4 ) : 8 in
let f = (lambda x : Real . (0.2 $ false : 64 / x)) :: Real -> {Bool : 0.2, Real : 0.8} in
(f n)

>>>

{
	 [Bool] false :: Bool : SymProb362 ,
	 [Real] 32. :: Real : SymProb363 ,
	 [Bool] false :: Bool : SymProb364 ,
	 [Real] 16. :: Real : SymProb365 ,
	 [Bool] false :: Bool : SymProb366 ,
	 [Real] 8. :: Real : SymProb367 
}  
\end{gplcblock}



\subsubsection{Condicional}

El bloque condicional en este lenguaje tiene
características similares a las que presenta
en otros lenguajes: El tipo de la condición
debe tener un tipo consistente con \lstinline|Bool|, y
los tipos de las ramas deben ser iguales, a
diferencia de la heterogeneidad que pueden tener
las ramas de la seleción binaria probabilística.
En \ref{lst:if1} el programa, la condición \lstinline|false| retorna \lstinline|-1|,
y ambas ramas tienen el mismo tipo.


\begin{gplcblock}[caption={if1},label={lst:if1}]
if false then 1 else (-1)

>>>

{
	 [Real] -1. :: Real : SymProb11 
}  
\end{gplcblock}

En cambio, el programa \ref{lst:if2} es rechazado
estáticamente debido a que las ramas de
ejecución tienen tipos diferentes.


\begin{gplcblock}[caption={if2},label={lst:if2}]
if true then 1 else true

>>>

Static Error [line 1, cols 0-20]: In (if), branches have unequal types {Real:1} and {Bool:1}
\end{gplcblock}


Los programas \ref{lst:if3} y \ref{lst:if4} comprueban que en
tiempo de ejecución no se ejecutan ambas ramas.
En el programa \ref{lst:if3}, se retorna exitosamente la
rama \lstinline|true|. En la rama false, el argumento de \lstinline|f|
pasa el chequeo estático de consistencia, porque
\lstinline|Bool| es consistente con el tipo \lstinline|?| del parámetro
formal. Sin embargo en el programa \ref{lst:if4}, la rama
false se ejecuta, y en tiempo de ejecución, el
programa es rechazado dinámicamente.

\begin{gplcblock}[caption={if3},label={lst:if3}]
let double = lambda x : ? . (x + x) in
if true then (double 1) else (double true) 

>>>

{
	 [Real] 2. :: Real : SymProb59 
}  
\end{gplcblock}


\begin{gplcblock}[caption={if4},label={lst:if4}]
let double = lambda x : ? . (x + x) in
if false then (double 1) else (double true)

>>>

Run-Time Error [line 1, cols 25-26]: expected value Real; found Bool
\end{gplcblock}


\begin{gplcblock}[caption={if5},label={lst:if5}]
let condition = (0.7 $ true : false) in
if condition then 100 else 200

>>>

{
	 [Real] 100. :: Real : SymProb59 ,
	 [Real] 200. :: Real : SymProb60 
}  
\end{gplcblock}


  
  Este programa demuestra dos aspectos relevantes
  del lenguaje y su implementación. En primer lugar,
la condición del bloque \lstinline|if| puede ser una variable
que recorre los elementos de una distribución. En
segundo lugar, el ejemplo muestra una forma efectiva
de implementar una selección binaria que fuerza la
homogeneidad de las ramas. En general,
\lstinline|let cond = (p $ true $ false) in (if cond then m else n)|
retorna un valor con un solo tipo en la distribución.


  
\subsubsection{Ascripción Simple}
La operación de ascripción con un tipo simple tiene
las siguientes funciones: Verificar estáticamente
la consistencia del tipo inferido de la expresión
ascrita con el tipo ascrito, y \lstinline|2|. Realizar un "cast"
del tipo de la expresión ascrita al tipo ascrito, ya
que estáticamente, el tipo de la expresión \lstinline|<expr>::<type>|
es el tipo simple ascrito \lstinline|<type>|.


El programa \ref{lst:ascr1} termina exitosamente,
la distribución de dirac con 1.0. Por otro lado
el programa \ref{lst:ascr2} falla estáticamente, porque
la expresión tiene tipo \lstinline|Bool|, inconsistente con
\lstinline|Real|. Finalmente, en el programa \ref{lst:ascr3}, el valor
\lstinline|true|, ascrito con \lstinline|?| se asigna a la variable
\lstinline|unknown|.  Esta ascripción funciona efectivamente
como un cast, que deja la variable con el tipo \lstinline|?|.
Por esta razón, el chequeo estático \lstinline|unknown :: Real|
no falla, porque \lstinline|?| es consistente con \lstinline|Real|.
Sin embargo, cuando se realiza el chequeo dinámico de
la ascripción en tiempo de ejecuión, la ascripción
falla debido a la incompatibilidad del valor con el
tipo Real.

\begin{gplcblock}[caption={ascr1},label={lst:ascr1}]
1 :: Real

>>>

{
	 [Real] 1. :: Real : SymProb1 
}  

\end{gplcblock}


\begin{gplcblock}[caption={ascr2},label={lst:ascr2}]
true :: Real

>>>

Static Error [line 1, cols 0-12]: In simple ascription, type Bool of term inconsistent with ascription type Real
\end{gplcblock}

\begin{gplcblock}[caption={ascr3},label={lst:ascr3}]
let unknown = (true :: ?) in
(unknown :: Real)

>>>

Run-Time Error [line 2, cols 1-8]: expected value Real; found Bool
\end{gplcblock}


En los programas \ref{lst:ascr4} y \ref{lst:ascr5} se demuestra
de otra forma el poder de la ascripción para
modificar el tipo de una expresión. En el programa
\ref{lst:ascr4}, el tipo del parámetro es \lstinline|?|. Dada la
consistencia del tipo \lstinline|Bool| con \lstinline|?|, \lstinline|(id false)|
se ejecuta sin problemas. Sin embargo, en el
programa \ref{lst:ascr5}
el tipo de \lstinline|id| se "refina" ascribiendo el tipo
de función \lstinline|Real -> {?:?}|, con dominio \lstinline|Real|.
La aplicación \lstinline|(real_id false)| falla estáticamente
por la inconsistencia del argumento con \lstinline|Real|,
demostrando el refinamiento del tipo esperado
para el argumento.


\begin{gplcblock}[caption={ascr4},label={lst:ascr4}]
let id = lambda x : ? . x in
(id false) 

>>>

{
	 [Bool] false :: ? : SymProb31 
}  
\end{gplcblock}





\begin{gplcblock}[caption={ascr5},label={lst:ascr5}]
let id = lambda x . x in
let real_id = id :: Real -> {?:?} in
(real_id false)

>>>

Static Error [line 3, cols 0-15]: In function, expected argument of type Real, but found Bool
\end{gplcblock}




\subsubsection{Integridad Semántica Probabilística}

El interprete tiene la capacidad para juzgar semánticamente:
\textbf{1.} bien o mal formuladas las probabilidades graduales usadas
en la selección probabilística y en la construcción de tipos de
distribución, y \textbf{2.} bien o mal construidas las distribuciones
de tipos usadas en la ascripción.

En el caso de las probabilidades graduales, la probabilidad no
especificada \lstinline|?| es siempre correctamente utilizada. La
probabilidad numérica, en cambio, debe ser un número real en el intervalo $[0,1]$. En el código \ref{lst:perr1}, la selección es rechazada estáticamente
debido al valor de la probabilidad.

\begin{gplcblock}[caption={ascr5},label={lst:perr1}]
1.1 $ true : false

>>>

Static Error [line 1, cols 0-3]: Expected probability in [0,1] but found 1.1
\end{gplcblock}


En el caso de las distribuciones, la regla con respecto a la
integridad de las probabilidades tiene 3 casos: \textbf{(1)} Si la
distribución solo contiene probabilidades numéricas, estas deben sumar
exactamente 1.0. \textbf{(2)} Si solo hay probabilidades desconocidas,
la distribución es correcta. \textbf{(3)} Si la distribución combina
probabilidades numéricas y no especificadas, entonces la suma de las
probabilidades numéricas debe ser un número en el intervalo $[0,1]$.
La distribución del código \ref{lst:perr2} es rechazado
estáticamente porque sus probabilidades numéricas suman más de 1.0.
La distribución del código \ref{lst:perr3} falla con el mismo mensaje, pero esta vez debido a que sus probabilidades
numéricas suman menos de 1.0, y no hay ninguna probabilidad \lstinline|?| que pueda complementar la probabilidad restante

\begin{gplcblock}[caption={Distribución construida },label={lst:perr2}]
(? $ true : 23) :: {Bool:0.5, Real:0.6, Real : ?}

>>>

Static Error [line 1, cols 19-49]: Invalid distribution type. The sum of probabilities must be 1.0.
If the sum of numeric probabilities is in [0, 1], the distribution must contain at least one '?' to complete it.
\end{gplcblock}

\begin{gplcblock}[caption={Distribución con Probabilidad},label={lst:perr3}]
(? $ true : 23) :: {Bool:0.4, Real:0.4}
\end{gplcblock}


\subsection{Resumen}

En conjunto, los casos presentados entregan evidencia práctica de que
el intérprete se comporta de manera consistente frente a los
escenarios estudiados: acepta programas bien formados y produce
resultados con la estructura esperada (distribuciones, sus valores y
las ascripciones asociadas), y rechaza programas inválidos cuando se
violan restricciones del lenguaje, reportando diagnósticos precisos
que incluyen localización y explicación de la inconsistencia
detectada. Además, los últimos ejemplos muestran que las reglas
semánticas sobre probabílidades válidas y distribuciones bien
construídas son efectivamente forzadas por el sistema.

Si bien esta validación no pretende ser exhaustiva, sí es sistemática
en el sentido de cubrir cada construcción central del lenguaje y una
colección de escenarios representativos y de borde. Por lo tanto, el
capítulo funciona tanto como una validación funcional del intérprete
como una guía concreta de uso del lenguaje GPLC, conectando
directamente las decisiones de diseño y las reglas semánticas con el
comportamiento observable de la implementación.
